%% ---------------------------------------------------------------------
\section{非微扰地位与可解模型}\label{sec:power-np}

领头幂次之外的一切，最终都依赖于我们对 QCD 非微扰结构的了解。本节先集中说明全书的非微扰地位，再介绍三个能真正非微扰地检验 OPE 与部分子图像的“实验室”。

\begin{remark}[非微扰地位]\label{rem:power-np-status}
\begin{enumerate}[label=(\roman*),leftmargin=1.8em,itemsep=1pt,topsep=2pt]
\item 本书的每个因子化定理，要么是重整化微扰论中关于红外正规化的部分子态的陈述（\tagR，$R_{\mathrm{PT}}$，或 \tagA），要么是关于微扰论加编号假设 \cref{asm:intro-wightman}、\cref{asm:intro-composite}、\cref{asm:intro-np-ope}、\cref{asm:intro-af-coeff}、\cref{asm:intro-growth}、\cref{asm:intro-bridge} 的陈述（\cref{def:intro-MA}）。QCD 尚未作为四维量子场论被构造（四维 Yang--Mills 理论的存在性与质量间隙是千禧年问题之一），所以“对强子严格”永远意味着“在所列非微扰假设之下证明完整”（$R_{\mathrm{cond}}$），[R] 本身从不意味着对真实世界 QCD 的数学证明。% M28
\item 在\cref{asm:intro-wightman}之下有真正的非微扰定理：谱函数的解析性与正定性（\cref{prop:ee-analytic}）、有限减除色散关系（\cref{prop:dis-dispersion-np}），以及另加非微扰 OPE 的 DIS 矩定理（\cref{thm:dis-moments-np}）；部分子分布的支撑（\cref{prop:pdf-support-hadron}，\tagA）与求和规则（\cref{prop:pdf-sum-rules}）也属于这一层次。最干净的非微扰连接是定域扭度 2 矩阵元的格点计算（第~22~章的\cref{prop:lattice-moments}）：它提供矩阵元，不检验强子截面的因子化。
\item 幂次压低的逐阶含义：每一阶的余项在积分范数下被 $O(\lambda^aL^K)$ 控制并对维数正规子一致（\cref{def:method-power-suppressed}），常数随阶数增长，这些界不能对 $\as$ 求和；展开是渐近的，而重整化子表明微扰级数与幂次修正不能被干净地分开。
\item 一般方法给出的指数可以远非最优。单圈离壳三角图中，一般估计（\cref{lem:appN-error-table}）在软区只保证相对 $O(\lambda^{1/9})$ 的误差，而精确余项是完整的一个幂次、次领头幂次无对数：$\int r=-\frac{4m^2}{Q^4}+\frac{m^4}{Q^6}(L^2-15+\frac{\pi^2}{3})+O(\frac{m^6}{Q^8}L^2)$，$L=\ln(Q^2/m^2)$（\cref{prop:appN-triangle-exact}(d)）；在纤维化积分范数下余项被 $C\lambda^2L^2$ 控制（\cref{prop:appN-triangle-cs}）。
\item 公理化场论中，OPE 可在一个相空间条件下被非微扰地导出\cite{Bostelmann:2005rr}（据该文摘要），该条件对 QCD 尚未建立。
\item 非微扰效应可以对 $\as$ 的任何幂次都不可见：瞬子贡献 $\propto\ee^{-2\pi/\as(Q)}=(\Lambda/Q)^{\beta_0}$（单圈），对偶性破坏在 OPE 的幂次展开中不可见（\cref{conj:ee-duality}）；二者与\cref{asm:intro-af-coeff}作为渐近陈述相容。
\item 可解模型给出相反的教训：二维 $U(N)$ Gross--Neveu 模型中扭度展开在 $1/N$ 的领头阶绝对收敛，而超可重整的二维 $O(N)$ 模型中动量空间的幂次展开因阶乘增强而发散、坐标空间的展开收敛（第~23~章的\cref{thm:yzliu-gn-convergence,prop:yzliu-on-factorial}）。它们是模型结果，不证明 QCD 中扭度展开的收敛或发散。
\end{enumerate}
\end{remark}

\subsection{微扰 \texorpdfstring{$\varphi^4_4$}{phi4} 中 OPE 的收敛}

\begin{theorem}[微扰 $\varphi^4_4$ 中 OPE 的收敛性 \tagR]\label{thm:power-ope-convergence}
（$R_{\mathrm{PT}}$：在重整化微扰论中对红外正规化的部分子态严格。）（本定理只涉及 $\varphi^4_4$ 模型的微扰论。）考虑四维欧氏空间中质量 $m>0$ 的厄米标量场 $g\varphi^4$ 理论，取定圈数 $l$。复合场 $O_A=\partial^{w_1}\varphi\cdots\partial^{w_n}\varphi$ 的量纲为 $[A]=n+\sum_i|w_i|$；旁观场为 $\varphi(f_{p_i})=\int\dd^4x\,\varphi(x)f_{p_i}(x)$，$f_{p_i}$ 光滑，$\hat f_{p_i}$ 的支撑在以 $p_i$ 为心、半径 $\epsilon$ 的球内。把 OPE $O_A(x)O_B(0)\sim\sum_C\mathcal C^C_{AB}(x)O_C(0)$ 的求和限于 $[C]-[A]-[B]\le\Delta$，则在 $l$ 圈阶有
\begin{multline}\label{eq:power-hk}
\Big|\Big\langle O_A(x)O_B(0)\prod_{i=1}^n\varphi(f_{p_i})\Big\rangle-\sum_C\mathcal C^C_{AB}(x)\Big\langle O_C(0)\prod_{i=1}^n\varphi(f_{p_i})\Big\rangle\Big|
\le m^{[A]+[B]+n}\sqrt{[A]!\,[B]!}\,\tilde K^{[A]+[B]}\prod_i\sup|\hat f_{p_i}|\\
\times\Big(\sup\big(1,\tfrac{|\vec p|_n}{m}\big)\Big)^{2([A]+[B])(n+2l+1)+3n}\sum_{\lambda=0}^{n/2+2l}\frac{\ln^\lambda\sup(1,\frac{|\vec p|_n}{m})}{2^\lambda\lambda!}
\cdot\frac{1}{\sqrt{\Delta!}}\Big(\tilde K\,m|x|\sup\big(1,\tfrac{|\vec p|_n}{m}\big)^{n+2l+1}\Big)^{\Delta},
\end{multline}
其中 $\tilde K$ 只依赖于 $n$ 与 $l$，$|\vec p|_n\equiv\sup_{J\subset\{1,\dots,n\}}|\sum_{i\in J}p_i|$\cite[定理~3]{Hollands:2011gf}。由于对任意 $c$ 有 $c^\Delta/\sqrt{\Delta!}\to0$，在每一圈阶上 OPE 在每个固定的 $|x|$ 处收敛，而不只是渐近成立。
% VF-FRO-4
% uses: none
\end{theorem}

\begin{proofsketch}
证明\cite{Hollands:2011gf}采用 Wilson--Wegner--Polchinski 重整化群流方程：引入红外截断与紫外截断 $\Lambda_0$，对带复合算符插入的连通截肢 Green 函数，在圈数与外腿数上归纳地给出对 $\Lambda_0$ 一致的界；关键是一个能在归纳中自我复制、同时控制 OPE 余项的归纳假设。完整证明见\cite[定理~3 及其证明]{Hollands:2011gf}；无质量情形见\cite{Holland:2014pna}，弯曲时空中微扰 OPE 的构造见\cite{Hollands:2006ag}（二者均据摘要）。
\end{proofsketch}

收敛性依赖于旁观场所代表的“态”（$\hat f_{p_i}$ 的支撑越不紧，收敛越差），不涉及圈数展开本身的收敛，在闵氏空间中只能期望对类空分离成立\cite{Hollands:2011gf}。这是模型陈述：对 QCD，微扰 OPE 是全阶论证（\cref{thm:dis-ope}，\tagA），其收敛性是猜想；更多严格的 OPE 结果（纯 Yang--Mills、共形场论中的光线 OPE 等）见第~24~章（\cref{prop:front-ope-scalar,prop:front-ope-ym,conj:front-ope-convergence,rem:front-ope-qcd-status}）。

\subsection{二维大 \texorpdfstring{$N_c$}{Nc} 的 't Hooft 模型}

二维 QCD 在 $N_c\to\infty$ 时可解\cite{tHooft1974}：夸克被禁闭、介子谱离散，深度非弹散射却仍可精确计算。取光锥坐标 $x^\pm=(x^0\pm x^1)/\sqrt2$ 与规范 $A^+=A_-=0$，则 $G^a_{+-}=-\partial_-A^a_+$，$A_+$ 没有传播自由度，其“传播子” $\ii\delta^{ab}/(k^+)^2$ 与 $k^-$ 无关（$k^+=0$ 处以截断 $|k^+|>\lambda$ 理解），顶点为 $\ii g\gamma^+t^a$；在 $N_c\to\infty$、$m_0^2\equiv g^2N_c/(2\pi)$ 固定时只剩彩虹自能与梯形图。彩虹自能 $\Sigma=-\ii\frac{g^2N_c}{4\pi}\gamma^+(\frac1\lambda-\frac{1}{p^+})$ 只依赖于 $p^+$，故单胶子结果即精确解，夸克传播子分母为 $2p^+p^--(m^2-m_0^2)-m_0^2|p^+|/\lambda+\ii0$：$\lambda\to0$ 时夸克极点被推到无穷远。对夸克（味 $a$）—反夸克（味 $b$）的 Bethe--Salpeter 振幅求梯形和，令 $\phi(x)=\int\dd k^-\,\Psi(xP^+,k^-)$；核与 $k^-$ 无关，$k^-$ 积分由留数完成，两个极点只在 $0<x<1$ 时位于实轴两侧，故 $\phi$ 的支撑在 $[0,1]$ 内。与 $\lambda$ 有关的项相消后得到 't Hooft 方程\cite{tHooft1974}
\begin{equation}\label{eq:power-thooft-eq}
\mu^2\phi(x)=\Big(\frac{m_a^2-m_0^2}{x}+\frac{m_b^2-m_0^2}{1-x}\Big)\phi(x)-m_0^2\,\mathrm{f.p.}\!\!\int_0^1\!\dd y\,\frac{\phi(y)}{(y-x)^2},\qquad\mu^2=2P^+P^-,
\end{equation}
f.p.\ 为 Hadamard 有限部分。由 $\mathrm{f.p.}\int_0^1\dd y\,(y-x)^{-2}=-\frac1x-\frac1{1-x}$，它等价于下面第一式，对应的二次型为第二式：
\begin{gather}
\mu^2\phi(x)=\Big(\frac{m_a^2}{x}+\frac{m_b^2}{1-x}\Big)\phi(x)-m_0^2\,\mathrm{P}\!\!\int_0^1\!\dd y\,\frac{\phi(y)-\phi(x)}{(y-x)^2},\label{eq:power-thooft-pv}\\
\langle\phi,H\phi\rangle=\int_0^1\!\dd x\Big(\frac{m_a^2}{x}+\frac{m_b^2}{1-x}\Big)|\phi(x)|^2+\frac{m_0^2}{2}\int_0^1\!\!\int_0^1\!\dd x\,\dd y\,\frac{|\phi(x)-\phi(y)|^2}{(x-y)^2}\;\ge0 .\label{eq:power-thooft-form}
\end{gather}

\begin{proposition}['t Hooft 模型 \tagP]\label{prop:power-thooft-model}
（在二维 QCD 的 $N_c\to\infty$ 极限中精确；作为关于 QCD 的陈述只是部分的。）在上述模型中：(a) 介子的质量与光锥波函数由\eqref{eq:power-thooft-eq}的本征值问题给出；谱离散、非负，$m_a=m_b=0$ 时最低态无质量且 $\phi_0=1$；$x\to0$ 时 $\phi_n\sim x^{\beta_a}$，$\pi\beta_a\cot(\pi\beta_a)=1-m_a^2/m_0^2$（$x\to1$ 处同理）；高激发态满足 $\mu_n^2\simeq\pi^2m_0^2\,n$。(b) 介子 $n$ 的夸克与反夸克分布为 $q_n(x)=|\phi_n(x)|^2$、$\bar q_n(x)=|\phi_n(1-x)|^2$，支撑在 $[0,1]$ 内，满足 $\int_0^1q_n=1$ 与 $\int_0^1x\,[q_n+\bar q_n]\,\dd x=1$（胶子不携带动量）。(c) 在 Bjorken 极限下，介子的深度非弹结构函数标度化，并且（对末态的离散共振取平均，即作为 $x$ 的分布弱收敛）等于以 (b) 中分布表示的部分子模型结果，修正是 $m_0^2/Q^2$、$m_a^2/Q^2$ 的幂次，没有对数型的标度破坏\cite{Einhorn1976}。这是规范理论中从第一原理非微扰地算出束缚态部分子分布的一个例子；该模型中的准部分子分布见第~23~章（\cref{prop:yzliu-2dqcd}）。
（本命题只涉及二维大 $N_c$ 模型。）
\uses{def:pdf-quark}
\end{proposition}

\begin{proof}
(a) 等价性由上述有限部分恒等式给出；\eqref{eq:power-thooft-form}由\eqref{eq:power-thooft-pv}关于 $x,y$ 对称化得到，两项都非负，故 $\mu^2\ge0$。$m_a=m_b=0$ 时第一项消失，第二项对常数为零，故 $\phi_0=1$、$\mu_0^2=0$（手征极限下无质量的“$\pi$ 介子”）。把 $\phi\simeq x^\beta$ 代入\eqref{eq:power-thooft-eq}并保留 $x\to0$ 时最奇异的项，用 $\mathrm{f.p.}\int_0^\infty\dd y\,y^\beta(y-x)^{-2}=-\pi\beta\cot(\pi\beta)\,x^{\beta-1}$（$-1<\beta<1$），得 $(m_a^2-m_0^2)+m_0^2\pi\beta\cot(\pi\beta)=0$。对高度振荡的 $\phi_n\simeq\sqrt2\sin(n\pi x)$，主值核作用于 $\ee^{\ii ky}$ 给出 $\pi|k|\ee^{\ii kx}$，故 $\mu_n^2\simeq\pi^2m_0^2n$。$H$ 是 $L^2[0,1]$ 上预解式为紧的正自伴算符，故谱离散、本征函数完备\cite{tHooft1974,Einhorn1976}。
(b) 光锥规范中 Wilson 线为 1，部分子分布（\cref{def:pdf-quark}的二维形式）是 $x^+=0$ 上数密度算符的矩阵元；$N_c\to\infty$ 时介子价夸克—反夸克分量的波函数就是 $\phi_n$，更高的 Fock 分量被 $1/N_c$ 压低\cite{Einhorn1976}，故夸克动量分数的几率密度为 $|\phi_n(x)|^2$，反夸克携带 $1-x$；归一化给出数目求和规则，$x+(1-x)=1$ 给出动量求和规则，支撑来自 (a)。
(c) Bjorken 极限下被击夸克在长度为 $O(1/q^-)$ 的光锥时间间隔内传播；耦合 $g^2$ 有质量平方的量纲，每次瞬时相互作用的插入都带来 $m_0^2/Q^2$ 量级的压低，故结构函数等于“手提包”项，即 (b) 中分布的部分子模型表达式；二维中没有对数发散的重整化，因而没有对数。固定 $Q^2$ 时结构函数是 $x$ 的一串 $\delta$ 函数，相邻共振的间距按 $m_0^2/Q^2$ 趋于零，所以标度化在分布意义下成立。对全部末态求和精确再现部分子模型，已由 Einhorn 明确验证（同时验证了 Drell--Yan--West 关系与 Bloom--Gilman 对偶）\cite{Einhorn1976}；电流振幅的构造与幺正性见\cite{CallanCooteGross1976}。（上述积分、等价形式、正定性、零模、端点指数与谱的渐近行为都用符号与数值计算核对过。）
\end{proof}

\paragraph{地位与缺口}
作为关于 QCD 的陈述这是 \tagP：模型在二维（无横向自由度、胶子不传播、超可重整，因而没有对数型标度破坏、Sudakov 对数与 Glauber 问题）并取 $N_c\to\infty$；谱的完备性依赖于所引的分析，(c) 中对末态共振的求和由\cite{Einhorn1976}完成。它说明禁闭本身不破坏部分子模型意义下的因子化、束缚态的部分子分布可以从第一原理算出；它不说明四维 QCD 中的任何东西（\cref{asm:intro-bridge}的地位不因此改变）。

\paragraph{第三个“实验室”：格点。}
格点 QCD 提供非微扰的数值证据（不是证明）：它计算定域算符的矩阵元，如扭度 2 算符的矩（第~22~章），也检验重整化子的预言（\cref{prop:power-renormalon-examples}(g)）；它不检验强子截面的因子化本身，% critic C12
其从欧氏关联函数到光锥分布的匹配又依赖于自己的因子化定理（第~22~章）。

\begin{table}[htbp]
\centering\small
\caption{本章陈述的地位一览。“性质”一栏说明陈述成立的意义。}
\label{tab:power-status}
\begin{tabular}{@{}p{3.3cm}p{3.4cm}p{1.6cm}p{2.2cm}p{3.5cm}@{}}
\toprule
陈述 & 标签 & 等级 & 性质 & 主要缺口 \\
\midrule
DIS 扭度 4（$x$ 空间） & \cref{prop:power-dis-twist4-x} & \tagA & E-S & 次领头幂次区域分析；重整化子混合 \\
SIDIS 型扭度 3 SSA & \cref{prop:power-twist3-sidis} & \tagP & 有限阶 & 无全阶论证 \\
$pp$ 中扭度 3 SSA & \cref{conj:power-twist3-pp} & \tagC & --- & 带色末态 \\
次软定理 & \cref{prop:power-next-to-soft} & \tagR & 树图，精确 & 圈图层次为 \tagA \\
NLP 端点发散 & \cref{prop:power-nlp-endpoint} & \tagP & 逐个过程 & 无一般定理 \\
OPE 基的重排 & \cref{prop:power-rearrangement} & \tagR & 代数恒等式 & --- \\
大 $\beta_0$ Adler 函数 & \cref{prop:power-adler-large-beta0} & \tagP & 极限中精确 & 朴素非阿贝尔化 \\
重整化子的例子；模型 & \cref{prop:power-renormalon-examples} & \tagP & 近似／模型 & 大 $\beta_0$ 之外无控制 \\
重整化子 $\leftrightarrow$ 高扭度 & \cref{conj:power-renormalon-ht} & \tagC & --- & 无非微扰 OPE \\
$\varphi^4_4$ 中 OPE 收敛 & \cref{thm:power-ope-convergence} & \tagR & 每阶显式界 & 模型；QCD 中为猜想 \\
't Hooft 模型 & \cref{prop:power-thooft-model} & \tagP & 模型中精确 & 二维，$N_c\to\infty$ \\
\bottomrule
\end{tabular}
\end{table}

\paragraph{本章小结。}
领头幂次之外，严格性迅速下降（\cref{tab:power-status}）：可以完整证明的是代数性与模型中的陈述；DIS 扭度 4 的 $x$ 空间形式有全阶论证；扭度 3 单自旋不对称与次领头幂次因子化只在有限阶与特定过程中得到验证；重整化子与高扭度的一般对应及强子—强子碰撞中的扭度 3 因子化是猜想。贯穿全章的一点是：超出领头幂次后，“微扰系数”与“非微扰矩阵元”的划分本身就有与幂次修正同量级的歧义（\cref{prop:power-rearrangement}，\cref{ex:power-renormalon-chain}），因此“精确到幂次修正”只有在幂次修正项被定义之后才有确切含义——这正是\cref{asm:intro-bridge}(Br3) 所表达的内容。
